\documentclass[11pt]{article}

\usepackage{amsmath,amssymb,amsthm,mathtools}
\usepackage{geometry}
\usepackage{hyperref}
\usepackage{enumitem}

\geometry{margin=1in}

\title{Quantization as a Change of State Space:\\
A Geometric and Convex-Structural Reformulation}
\author{}
\date{}

\newtheorem{theorem}{Theorem}
\newtheorem{definition}{Definition}
\newtheorem{proposition}{Proposition}
\newtheorem{remark}{Remark}

\begin{document}

\maketitle

\begin{abstract}
Canonical quantization is traditionally formulated as a map from classical observables,
equipped with a Poisson bracket, to noncommuting operators satisfying commutation relations.
However, it is known that no such map can preserve all desired structural properties.
In this article we reformulate quantization at a deeper level:
not as a homomorphism of observable algebras,
but as a deformation of the classical state space.
We show that the essential structural transition from classical to quantum theory
is the replacement of a simplex of probability measures by a non-simplex convex state space.
In this framework, the obstruction to Dirac's canonical quantization program
is understood geometrically as the impossibility of preserving the convex geometry
of states while deforming only the observable algebra.
\end{abstract}

\section{Introduction}

Canonical quantization is often described by a rule
\begin{equation}
\{f,g\}_{\mathrm{PB}} \longmapsto \frac{1}{i\hbar} [\hat f,\hat g],
\end{equation}
suggesting a map from classical observables to operators.
However, it is known that no map from the full Poisson algebra of smooth functions
to operators on a Hilbert space can preserve
both the Lie structure and the associative structure globally.

Rather than viewing this as a technical obstruction,
we argue that the difficulty is structural:
quantization does not primarily deform the algebra of observables.
It changes the geometry of the state space.

\section{Classical Theory as Commutative Probability}

Let $M$ be a classical phase space.

\subsection*{Observables}
\[
\mathcal{A}_{\mathrm{cl}} = C^\infty(M)
\]
with pointwise multiplication and Poisson bracket.

\subsection*{States}
States are probability measures:
\[
\mathcal{S}_{\mathrm{cl}} = \mathcal{P}(M).
\]

Expectation values are given by
\begin{equation}
\langle f \rangle_\mu = \int_M f \, d\mu.
\end{equation}

\begin{proposition}
The classical state space $\mathcal{P}(M)$ is a Choquet simplex.
\end{proposition}

\begin{proof}[Sketch]
Extreme points are Dirac measures $\delta_x$.
Every probability measure admits a unique barycentric decomposition
in terms of these extremal states.
\end{proof}

Thus classical mechanics is fundamentally
\emph{commutative probability theory}.
The algebra of observables is dual to the simplex of states.

\section{Quantum Theory as Noncommutative Probability}

Let $\mathcal{H}$ be a Hilbert space.

\subsection*{Observables}
Selfadjoint operators in $\mathcal{B}(\mathcal{H})$.

\subsection*{States}
Density operators:
\[
\mathcal{S}_{\mathrm{qm}} = 
\{ \rho \ge 0 \mid \mathrm{Tr}(\rho)=1 \}.
\]

Expectation values:
\begin{equation}
\langle A \rangle_\rho = \mathrm{Tr}(\rho A).
\end{equation}

\begin{proposition}
The quantum state space $\mathcal{S}_{\mathrm{qm}}$ is convex but not a simplex.
\end{proposition}

\begin{proof}[Sketch]
Pure states are rank-one projectors.
Mixed states admit infinitely many distinct decompositions into pure states.
Hence barycentric decomposition is not unique.
\end{proof}

This non-simplex structure is the essential geometric difference
between classical and quantum theory.

\section{Observables as Affine Functionals}

Observables can be reconstructed from states:

\begin{equation}
A \quad \longleftrightarrow \quad
\rho \mapsto \mathrm{Tr}(\rho A).
\end{equation}

Thus observables are affine functionals on the convex state space.

Therefore:

\begin{center}
Changing the state space automatically changes the observable algebra.
\end{center}

One cannot deform one independently of the other.

\section{Quantization as State-Space Deformation}

We now state the central structural principle.

\begin{definition}
Quantization is the replacement
\[
\mathcal{S}_{\mathrm{cl}} = \mathcal{P}(M)
\quad \longrightarrow \quad
\mathcal{S}_{\mathrm{qm}} = \mathcal{D}(\mathcal{H}),
\]
where a simplex state space is replaced by a non-simplex convex state space.
\end{definition}

Under this reformulation:

\begin{itemize}[leftmargin=1.5em]
\item Classical pure states are points $x \in M$.
\item Quantum pure states are rays in $\mathcal{H}$.
\item Classical observables are functions on $M$.
\item Quantum observables are affine functionals on $\mathcal{S}_{\mathrm{qm}}$.
\end{itemize}

No homomorphism between observable algebras is required.

\section{Geometric Interpretation of the Quantization Obstruction}

Suppose one attempts to construct a map
\[
Q: C^\infty(M) \to \mathcal{B}(\mathcal{H})
\]
preserving:

\begin{enumerate}
\item Poisson bracket $\to$ commutator,
\item Unit element,
\item Canonical covariance.
\end{enumerate}

Such a map would imply a dual affine embedding
of $\mathcal{S}_{\mathrm{qm}}$ into $\mathcal{P}(M)$
preserving convex structure.

\begin{theorem}[Geometric Obstruction]
There is no affine isomorphism between a simplex
and a non-simplex convex set.
\end{theorem}

\begin{proof}
Simplex property is invariant under affine isomorphism.
Quantum state space fails the uniqueness of decomposition property.
Therefore no such affine equivalence exists.
\end{proof}

Thus the obstruction to canonical quantization
is geometric rather than algebraic.

\section{Loss of Global Event Space}

Classical theory admits a single global sample space $M$.
All observables are jointly measurable functions on this space.

Quantum theory admits only contextual outcome spaces,
associated with commuting sets of observables.
No single global event space exists.

Quantization therefore represents:

\begin{center}
\textbf{the loss of a global underlying probability space.}
\end{center}

\section{Conclusion}

The traditional problem of quantization
asks for a homomorphism of observable algebras.
This problem has no solution in full generality.

However, when reformulated structurally,
quantization is not a map of algebras.
It is a change in convex geometry:

\begin{equation}
\text{Simplex state space}
\quad \longrightarrow \quad
\text{Non-simplex state space}.
\end{equation}

The noncommutativity of observables,
the failure of canonical covariance,
and contextuality
all follow naturally from this geometric transition.

Quantization is therefore best understood
not as a deformation of functions,
but as a deformation of probability itself.

\end{document}
